\section{GROUP ONE}\label{group-one}

This group was seen within paragraphing text (within a documentation
chunk).

\begin{verbatim}
{\Tt{}\nwlinkedidentq{whyse}{NW43JrUv-4YCtgL-1}\nwendquote}
\end{verbatim}

Reformatted, without arguments in the commands, it reads as follows.

\begin{verbatim}
{
    \Tt{}
    \nwlinkedidentq{}{}
    \nwendquote
}
\end{verbatim}

Whether all paragraphing text references are like this remains to be
seen.

Not all prosaic noweb macros are entirely composed of
\texttt{nwlinkedidentq}:

\begin{verbatim}
{\Tt{}M-x{\char126}customize-group{\char126}\nwlinkedidentq{whyse}{NW43JrUv-4YCtgL-1}\nwendquote}
\end{verbatim}

This time the reformatted group will not be stripped of command
arguments.

\begin{verbatim}
{
    \Tt{}
    M-x{\char126}customize-group{\char126}
    \nwlinkedidentq{whyse}{NW43JrUv-4YCtgL-1}
    \nwendquote
}
\end{verbatim}

The third and fourth groups (within the first) ensure that
\texttt{\textbackslash{}char126} is interpreted as a single unit,
inserting a character with that number.

\texttt{nwlinkedidentq} and \texttt{nwlinkedidentc} are commands which
call internal macros appropriate to the selected package options.

\textbf{Options applicable to all identifiers:} - subscriptidents; -
nosubscriptidents; - hyperidents; - nohyperidents;

\textbf{Options applicable only to identifiers appearing in quoted
code:} - subscriptquotedidents; - nosubscriptquotedidents; -
hyperquotedidents; - nohyperquotedidents.

\section{GROUP TWO}\label{group-two}

There command \texttt{nwlinkedidentc} is markup for identifiers used
\emph{in code}:

\begin{verbatim}
\nwfilename{/var/home/bryce/Documents/src/whyse/src/whyse.nw}\nwbegincode{1}\sublabel{NW43JrUv-4YCtgL-1}\nwmargintag{{\nwtagstyle{}\subpageref{NW43JrUv-4YCtgL-1}}}\moddef{Customization and global variables~{\nwtagstyle{}\subpageref{NW43JrUv-4YCtgL-1}}}\endmoddef\nwstartdeflinemarkup\nwusesondefline{\\{NW43JrUv-4e7Hxw-3}}\nwprevnextdefs{\relax}{NW43JrUv-4YCtgL-2}\nwenddeflinemarkup
\nwindexdefn{\nwixident{whyse}}{whyse}{NW43JrUv-4YCtgL-1}(defgroup \nwlinkedidentc{whyse}{NW43JrUv-4YCtgL-1} nil
  "noWeb HYpertext System in Emacs"
  :tag "WHYSE"
  :group 'applications)

\nwindexdefn{\nwixident{w-registered-projects}}{w-registered-projects}{NW43JrUv-4YCtgL-1}(defcustom \nwlinkedidentc{w-registered-projects}{NW43JrUv-4YCtgL-1} nil
  "This variable stores all of the projects that are known to WHYSE."
  :group 'whyse
  :type '(repeat w--project-widget)
  :require 'widget
  :tag "WHYSE Registered Projects")

\nwalsodefined{\\{NW43JrUv-4YCtgL-2}\\{NW43JrUv-4YCtgL-3}}\nwused{\\{NW43JrUv-4e7Hxw-3}}\nwidentdefs{\\{{\nwixident{w-registered-projects}}{w-registered-projects}}\\{{\nwixident{whyse}}{whyse}}}\nwendcode{}\nwbegindocs{2}\nwdocspar
\end{verbatim}

Reformatting this longer section, we have the following:

\begin{verbatim}
\nwfilename{/var/home/bryce/Documents/src/whyse/src/whyse.nw}
\nwbegincode{1}
\sublabel{NW43JrUv-4YCtgL-1}
\nwmargintag{
    {
        \nwtagstyle{}
        \subpageref{NW43JrUv-4YCtgL-1}
    }
}
\moddef{Customization and global variables~
    {
        \nwtagstyle{}
        \subpageref{NW43JrUv-4YCtgL-1}
    }
}
\endmoddef

\nwstartdeflinemarkup
    \nwusesondefline{\\{NW43JrUv-4e7Hxw-3}}
    \nwprevnextdefs{\relax}{NW43JrUv-4YCtgL-2}
\nwenddeflinemarkup
\nwindexdefn{\nwixident{whyse}}{whyse}{NW43JrUv-4YCtgL-1}% three arguments
(defgroup \nwlinkedidentc{whyse}{NW43JrUv-4YCtgL-1} nil
  "noWeb HYpertext System in Emacs"
  :tag "WHYSE"
  :group 'applications)

\nwindexdefn{\nwixident{w-registered-projects}}{w-registered-projects}{NW43JrUv-4YCtgL-1}
(defcustom \nwlinkedidentc{w-registered-projects}{NW43JrUv-4YCtgL-1} nil
  "This variable stores all of the projects that are known to WHYSE."
  :group 'whyse
  :type '(repeat w--project-widget)
  :require 'widget
  :tag "WHYSE Registered Projects")

\nwalsodefined{
    \\{NW43JrUv-4YCtgL-2}
    \\{NW43JrUv-4YCtgL-3}
}
\nwused{
    \\{NW43JrUv-4e7Hxw-3}
}
\nwidentdefs{
    \\{
        {\nwixident{w-registered-projects}}
        {w-registered-projects}
    }
    \\{
        {\nwixident{whyse}}{whyse}
    }
}
\nwendcode{}

\nwbegindocs{2}\nwdocspar
\end{verbatim}

\ldots and as formatted by \textbf{latexindent.pl}:

\begin{verbatim}
\nwfilename{/var/home/bryce/Documents/src/whyse/src/whyse.nw}
\nwbegincode{1}
\sublabel{NW43JrUv-4YCtgL-1}
\nwmargintag{
    {
        \nwtagstyle{}
        \subpageref{NW43JrUv-4YCtgL-1}
    }
}
\moddef{Customization and global variables~
    {
        \nwtagstyle{}
        \subpageref{NW43JrUv-4YCtgL-1}
    }
}
\endmoddef

\nwstartdeflinemarkup
\nwusesondefline{\\{NW43JrUv-4e7Hxw-3}}
\nwprevnextdefs{\relax}{NW43JrUv-4YCtgL-2}
\nwenddeflinemarkup
\nwindexdefn{\nwixident{whyse}}{whyse}{NW43JrUv-4YCtgL-1}% three arguments
(defgroup \nwlinkedidentc{whyse}{NW43JrUv-4YCtgL-1} nil
"noWeb HYpertext System in Emacs"
:tag "WHYSE"
:group 'applications)

\nwindexdefn{\nwixident{w-registered-projects}}{w-registered-projects}{NW43JrUv-4YCtgL-1}
(defcustom \nwlinkedidentc{w-registered-projects}{NW43JrUv-4YCtgL-1} nil
"This variable stores all of the projects that are known to WHYSE."
:group 'whyse
:type '(repeat w--project-widget)
:require 'widget
:tag "WHYSE Registered Projects")

\nwalsodefined{
    \\{NW43JrUv-4YCtgL-2}
\\{NW43JrUv-4YCtgL-3}
}
\nwused{
    \\{NW43JrUv-4e7Hxw-3}
}
\nwidentdefs{
    \\{
    {\nwixident{w-registered-projects}}
    {w-registered-projects}
}
\\{
{\nwixident{whyse}}{whyse}
}
}
\nwendcode{}

\nwbegindocs{2}\nwdocspar
\end{verbatim}

Until a later time, I can simply forget about the \texttt{nwfilename}
macro. I am working on the absolute basic support at the moment. I
cannot think about multi-noweb file documents yet.

\texttt{nwbegincode} and \texttt{nwendcode} are environment specifying
macros. It's fairly easy to understand that these must be paired
properly. Exactly what they do is not a concern to me yet, I'm just
noting that these are here and that they're mirrored commands. BEGIN.
END. Simple, right?

\texttt{moddef} and \texttt{endmoddef} are likewise paired commands.
(IMPLICITLY BEGIN)THING. ENDTHING. This is less simple.

\texttt{nwstartdeflinemarkup} and \texttt{nwenddeflinemarkup} are
likewise a pair.

\texttt{nwindexdefn} is a command that defines a term to be indexed. It
takes three arguments.

\texttt{nwlinkedidentc} is a cross-reference to an identifier. It takes
two arguments.
